\Needspace{14\baselineskip}
\section{Método}
\label{sec:metodo}

Esta revisão sistemática segue a metodologia de
\citeonline{kitchenham2007guidelines} para revisões sistemáticas em
Engenharia de Software, incorporando as atualizações de
\citeonline{kitchenham2013systematic}.
O reporte adota as Software Engineering Guidelines for Reporting
Secondary Studies (SEGRESS)~\cite{kitchenham2023segress}, adaptação do
PRISMA 2020~\cite{page2021prisma} para estudos secundários em Engenharia
de Software.
A estratégia de busca segue a extensão
PRISMA-S~\cite{rethlefsen2021prismas}, e o reporte completo segue o
checklist de 27 itens do PRISMA 2020.

\subsection{Protocolo e registro}
\label{sec:protocolo}

O protocolo foi elaborado antes do início das buscas.
Seis emendas foram registradas durante a execução, refinando critérios
de elegibilidade e o formulário de extração.

\subsection{Critérios de elegibilidade}
\label{sec:criterios}

A formulação dos critérios seguiu o \textit{framework} PICo:
a \textbf{P}opulação compreende agentes de codificação baseados em IA
com mecanismos de persistência de conhecimento entre sessões;
o \textbf{I}nteresse abrange as arquiteturas de persistência, sua
eficácia e custo-eficiência;
o \textbf{Co}ntexto inclui avaliações empíricas, benchmarks e propostas
de ferramentas ou frameworks.

A \autoref{tab:criterios} apresenta os critérios de inclusão e exclusão.
CI1 exige componente explícito de armazenamento e reutilização de
conhecimento entre sessões ou tarefas sequenciais;
sistemas que apenas recuperam contexto dentro de uma única sessão
não atendem ao critério.
CI2 estabelece piso mínimo de avaliação empírica com pelo menos duas
tarefas e métricas quantitativas.
CE2 opera como regra condicional: papers sobre agentes LLM genéricos
foram incluídos quando a avaliação continha benchmarks de Engenharia
de Software.
O período de cobertura abrange publicações de 2023 a março de 2026.

\begin{table}[htbp]
  \centering
  \caption{Critérios de elegibilidade.}
  \label{tab:criterios}
  \begin{tabularx}{\textwidth}{l l X}
    \toprule
    ID  & Tipo & Descrição \\
    \midrule
    CI1 & Inclusão & Propõe ou avalia mecanismo para persistir
      conhecimento entre sessões em agentes de codificação baseados
      em IA \\
    CI2 & Inclusão & Reporta resultados empíricos com pelo menos
      duas tarefas e métricas quantitativas \\
    CI3 & Inclusão & Publicado entre 2023 e março de 2026 \\
    CI4 & Inclusão & Escrito em inglês ou português \\
    CI5 & Inclusão & Publicação com revisão por pares ou preprint com
      autores e instituições identificáveis \\
    \midrule
    CE1 & Exclusão & Aborda apenas RAG ou recuperação de contexto
      intra-sessão, sem persistência entre sessões \\
    CE2 & Exclusão & Aborda apenas tarefas de linguagem natural, sem
      avaliação em benchmarks de Engenharia de Software \\
    CE3 & Exclusão & Avaliação limitada a exemplos anedóticos \\
    CE4 & Exclusão & Tutorial, editorial ou position paper sem
      contribuição técnica \\
    CE5 & Exclusão & Versão duplicada ou supersedida por versão posterior
      dos mesmos autores \\
    \bottomrule
  \end{tabularx}
\end{table}

\subsection{Fontes de informação}
\label{sec:fontes}

As buscas foram realizadas em duas bases de dados com suporte a busca
booleana reproduzível: Scopus (via API) e arXiv (via API REST com
paginação completa).
A estratégia de duas bases complementada por rastreamento de citações é
validada empiricamente por \citeonline{mourao2020hybrid}, que demonstraram
que Scopus combinado com \textit{snowballing} alcança \textit{recall} comparável a buscas em
quatro ou mais bases para revisões sistemáticas em Engenharia de Software.
O IEEE Xplore não foi buscado separadamente porque o Scopus indexa
integralmente seus proceedings e journals.

A busca foi suplementada por rastreamento de citações forward e backward
via Semantic Scholar a partir de sete artigos-semente selecionados por
representarem diversidade arquitetural e serem os mais citados no
subcampo~\cite{du2026memory,joshi2025swebenchcl,wang2026codemem,deng2026memcoder,wu2025gcc,wang2026improving,lindenbauer2025ctimrover}.
Um dos sete artigos-semente (CodeMEM) não foi capturado pelas strings de
busca por utilizar terminologia fora do vocabulário das strings;
sua adição manual ao \textit{pool} de triagem foi registrada como emenda ao
protocolo.

\subsection{Estratégia de busca}
\label{sec:busca}

A string de busca combina três blocos conceituais conectados por AND:
um bloco sobre agentes de IA para código (C1, 12 termos) com um segundo
sobre memória e persistência (C2, 13 termos), restritos ao domínio de
engenharia de software por um terceiro bloco (C3, 8 termos).
A string foi adaptada para a sintaxe de cada base: TITLE-ABS-KEY no
Scopus com filtro \texttt{PUBYEAR\,>\,2022}, e campos \texttt{abs:}/\texttt{ti:}
no arXiv com filtro de categorias cs.SE, cs.AI e cs.CL.\footnote{String
Scopus: \texttt{TITLE-ABS-KEY(("LLM agent" OR "AI coding agent" OR
"software engineering agent" OR "autonomous coding" OR "code agent" OR
"LLM-based software" OR "AI-assisted software" OR "coding assistant" OR
"language agent" OR "code assistant" OR "agentic coding" OR "SWE agent")
AND ("memory" OR "knowledge persistence" OR "experience" OR "context
management" OR "long-term memory" OR "episodic memory" OR "knowledge
reuse" OR "continual learning" OR "experience replay" OR
"memory-augmented" OR "persistent memory" OR "persistent context" OR
"cross-task") AND ("software engineering" OR "software development" OR
"code generation" OR "bug fix" OR "code repair" OR "repository" OR
"software maintenance" OR "issue resolution")) AND PUBYEAR > 2022}.
A string completa para o arXiv e o checklist PRISMA-S estão disponíveis
no material suplementar.}

Antes da execução, as strings passaram por auditoria contra oito SLRs
comparáveis, 16 sistemas conhecidos e dez formulações alternativas de
termos, resultando na adição de seis termos e cobertura de seis dos sete
artigos-semente.
O único artigo-semente não recuperado (CodeMEM) utiliza terminologia fora
do vocabulário das strings e foi adicionado manualmente.

A \autoref{tab:busca} apresenta os resultados por fonte.
As buscas foram executadas em 25 de março de 2026.
O rastreamento de citações identificou 186 registros únicos (8 forward,
178 backward).
Após consolidação e desduplicação, o \textit{pool} inicial contém 1.111 registros
para triagem.

\begin{table}[htbp]
  \centering
  \caption{Resultados das buscas por fonte.}
  \label{tab:busca}
  \begin{tabular}{l r}
    \toprule
    Fonte & Registros \\
    \midrule
    Scopus            &    39 \\
    arXiv             &   968 \\
    Citações forward  &     8 \\
    Citações backward &   178 \\
    Manual            &     1 \\
    \midrule
    Total bruto       & 1.194 \\
    Duplicatas removidas & $-83$ \\
    \textbf{Total para triagem} & \textbf{1.111} \\
    \bottomrule
  \end{tabular}
\end{table}

\subsection{Processo de seleção}
\label{sec:selecao}

A seleção ocorreu em duas fases.
Os 1.111 registros foram triados por título e resumo pelo
pesquisador, aplicando os critérios CI/CE.
Registros classificados como incertos avançaram para a fase seguinte.
Dos 1.111 triados, 1.060 foram excluídos e 51 avançaram para avaliação
por texto completo.

A leitura integral dos 51 artigos resultou em 23 exclusões.
Dos 23 excluídos, 12 foram removidos por CE2 (agentes LLM sem avaliação
em Engenharia de Software) e 11 pela combinação CI1/CE1 (recuperação
intra-sessão sem persistência entre sessões).
Três casos limítrofes foram aceitos com justificativa documentada.
Ao final, 28 estudos foram incluídos para extração de dados.

\subsection{Extração de dados}
\label{sec:extracao}

Os dados foram extraídos com formulário de 37 campos organizados em
sete categorias, com ênfase em arquitetura de persistência (8 campos)
e custo e eficiência (6 campos), que correspondem diretamente às
questões de pesquisa.
Dezesseis campos utilizam vocabulário controlado com valores predefinidos.
Campos sem dados disponíveis recebem NR (não reportado) quando o paper não
menciona o dado ou NA (não aplicável) quando o campo não se aplica ao
sistema, conforme convenção registrada como emenda ao protocolo.

Cada paper foi lido, seus dados registrados no formulário e os valores
verificados em seguida.
A validação incluiu verificação programática de 448 valores enum e
\textit{spot-check} de 40 campos interpretativos em dez papers.
A taxa de erro final foi inferior a 1\%, correspondente a oito
correções em 1.036 valores extraídos (28 papers $\times$ 37 campos).

\subsection{Avaliação de qualidade}
\label{sec:qualidade}

A qualidade dos estudos foi avaliada com checklist de seis itens
adaptado de \citeonline{kitchenham2007guidelines}: (Q1) objetivos
claramente definidos, (Q2) método de avaliação adequado, (Q3) ameaças
à validade discutidas, (Q4) resultados baseados em dados empíricos,
(Q5) contexto descrito para replicação e (Q6) métricas claramente
definidas.
Cada item recebeu pontuação 0, 0,5 ou 1, com máximo de 6,0.

A distribuição das pontuações varia de 2,0 a 6,0, com mediana de 5,0.
Vinte e seis dos 28 estudos (93\%) obtiveram pontuação igual ou
superior a 4,0.
Os dois estudos com pontuação inferior a
3,0~\cite{bui2026opendev,du2026memory} contribuem apenas descrições
qualitativas e não impactam os achados quantitativos, conforme
verificado na análise de sensibilidade (\autoref{sec:resultados}).
Para calibração, oito papers foram reavaliados após a extração completa.

\subsection{Métodos de síntese}
\label{sec:sintese}

A síntese é narrativa, organizada tematicamente pelas cinco questões
secundárias (QS1 a QS5).
Meta-análise foi descartada pela heterogeneidade de benchmarks, modelos
base e condições experimentais entre os estudos.

Duas análises de sensibilidade foram conduzidas:
(S1) re-síntese restrita aos cinco estudos publicados em venues com
revisão por pares, para verificar se os achados se sustentam sem
preprints;
(S2) exclusão dos dois estudos com pontuação inferior a 3,0, avaliando
o impacto nos achados quantitativos.

A avaliação de viés de publicação formal não é aplicável para síntese
narrativa~\cite{page2021prisma}.
O predomínio de preprints no corpus (82\%) é reconhecido como fonte
potencial de viés e discutido nas limitações.
A certeza da evidência foi avaliada qualitativamente por achado:
alta (múltiplos estudos de qualidade $\geq$~4/6 convergem),
moderada (poucos estudos ou qualidade variável) e
baixa (achado baseado em estudo único ou de baixa
qualidade)~\cite{kitchenham2023segress}.
